\documentclass[12pt,a4paper]{article}

\usepackage[utf8]{inputenc}
\usepackage[T1]{fontenc}
\usepackage[brazil]{babel}
\usepackage[top=3cm,left=3cm,bottom=2cm,right=2cm]{geometry}
\usepackage{mathptmx}
\usepackage{inconsolata}
\usepackage{setspace}
\usepackage{indentfirst}
\usepackage{titlesec}
\usepackage{booktabs}
\usepackage{array}
\usepackage{listings}
\usepackage{xcolor}
\usepackage{xurl}
\urlstyle{same}
\usepackage{amsmath}

\onehalfspacing
\setlength{\parindent}{1.25cm}

\titleformat*{\section}{\normalfont\bfseries\normalsize}
\titleformat*{\subsection}{\normalfont\bfseries\normalsize}

\lstset{
    basicstyle=\ttfamily\footnotesize,
    breaklines=true,
    columns=fullflexible,
    keepspaces=true,
    frame=single,
    rulecolor=\color{black!25},
    backgroundcolor=\color{black!4},
    xleftmargin=0.2cm,
    xrightmargin=0.2cm,
    framexleftmargin=4pt,
    framexrightmargin=4pt,
    framextopmargin=4pt,
    framexbottommargin=4pt,
    literate=
        {á}{{\'a}}1 {é}{{\'e}}1 {í}{{\'i}}1 {ó}{{\'o}}1 {ú}{{\'u}}1
        {Á}{{\'A}}1 {É}{{\'E}}1 {Í}{{\'I}}1 {Ó}{{\'O}}1 {Ú}{{\'U}}1
        {ã}{{\~a}}1 {õ}{{\~o}}1 {Ã}{{\~A}}1 {Õ}{{\~O}}1
        {â}{{\^a}}1 {ê}{{\^e}}1 {ô}{{\^o}}1
        {à}{{\`a}}1 {ç}{{\c{c}}}1 {Ç}{{\c{C}}}1
}

\begin{document}

\section{Questão 13: Como funciona a arquitetura Transformer}

\subsection{Da sequência de tokens ao contexto}

Na frase ``O banco aprovou o empréstimo'', a palavra ``banco'' designa uma instituição financeira. Em ``O banco da praça estava molhado'', a mesma palavra se refere a um assento. O Transformer pode construir representações diferentes para o mesmo token porque a atenção permite combinar informações de outras posições da sequência. Vaswani e colaboradores apresentaram a arquitetura em 2017 para tarefas de tradução [1].

Primeiro, um tokenizador divide o texto em unidades do vocabulário, que podem ser palavras, fragmentos de palavras ou pontuação. Cada unidade é convertida em um vetor de entrada, o \textit{embedding}. O modelo original soma a esse vetor uma codificação posicional senoidal. Assim, as operações de atenção recebem informação sobre o conteúdo e a posição. Isso ajuda a distinguir sequências com as mesmas palavras em ordens diferentes.

\subsection{Codificador e decodificador}

\begin{table}[htbp]
    \centering
    \begin{tabular}{>{\raggedright\arraybackslash}p{4.3cm} >{\raggedright\arraybackslash}p{10.5cm}}
        \toprule
        \textbf{Componente} & \textbf{Função na tradução} \\
        \midrule
        Codificador & Aplica autoatenção entre posições da entrada e produz um vetor contextual para cada token. \\
        \addlinespace
        Decodificador & Usa autoatenção causal no texto de saída e atenção cruzada sobre os vetores do codificador. A partir disso, calcula a distribuição do próximo token. \\
        \addlinespace
        Rede feed-forward & Transforma cada posição separadamente com as mesmas matrizes dentro da camada. A atenção faz a troca de informação entre posições. \\
        \addlinespace
        Conexões residuais e normalização & Na arquitetura original, cada subcamada soma sua transformação ao vetor de entrada e aplica normalização por camada. \\
        \bottomrule
    \end{tabular}
\end{table}

As camadas são empilhadas, de modo que uma camada recebe as representações já contextualizadas pela anterior. Na atenção cruzada, as consultas vêm do decodificador, enquanto as chaves e os valores vêm do codificador. Essa ligação permite ao decodificador usar partes relevantes do texto de origem em cada passo da tradução.

\subsection{Atenção por produto escalar}

Cada cabeça projeta os vetores de entrada em consultas $Q$, chaves $K$ e valores $V$. A consulta de uma posição é comparada às chaves das posições acessíveis. O produto escalar mede a compatibilidade. A divisão pela raiz da dimensão das chaves ajuda a controlar a escala desses produtos. A softmax transforma a linha de escores em pesos positivos que somam 1. A saída é a soma ponderada dos vetores de valor [1].

\[
\text{Atenção}(Q, K, V) = \mathrm{softmax}(QK^{\top} / \sqrt{d_k})\,V
\]

\subsection{Exemplo numérico de uma cabeça}

Considere apenas três posições de ``O banco aprovou o empréstimo'', com ``empréstimo'' tratado como um token para simplificar. Vamos atualizar o vetor da posição ``banco''. Os vetores da tabela são didáticos. Em um modelo treinado, eles vêm das projeções aprendidas. A consulta é $q = (2, 0)$ e as chaves têm dimensão 2.

\begin{table}[htbp]
    \centering
    \begin{tabular}{>{\raggedright\arraybackslash}p{2.7cm} >{\raggedright\arraybackslash}p{2.3cm} >{\raggedright\arraybackslash}p{2.2cm} >{\raggedright\arraybackslash}p{3.0cm} >{\raggedright\arraybackslash}p{2.8cm}}
        \toprule
        \textbf{Posição} & \textbf{Chave k} & \textbf{Valor v} & \textbf{Escore escalado} & \textbf{Peso softmax} \\
        \midrule
        empréstimo & (1, 0) & (1, 0) & 1,414 & 0,576 \\
        \addlinespace
        aprovou & (0,5 e 1) & (0, 1) & 0,707 & 0,284 \\
        \addlinespace
        banco & (0, 1) & (1, 1) & 0,000 & 0,140 \\
        \bottomrule
    \end{tabular}
\end{table}

Em ``empréstimo'', o produto escalar vale 2, e o escore escalado vale $2/\sqrt{2}$. Após a softmax, essa posição recebe o maior peso, cerca de 57,6\%. A soma ponderada gera o vetor $(0,716 \text{ e } 0,424)$ para a posição ``banco''. Esse vetor incorpora informação dos três valores, com maior contribuição ponderada da posição ``empréstimo'' em termos de peso.

A atenção com múltiplas cabeças repete esse cálculo com projeções diferentes, concatena os resultados e aplica uma projeção final. Uma cabeça pode aprender a dar peso a palavras relacionadas ao sentido de ``banco'', enquanto outra pode responder a um padrão de posição ou de sintaxe. Esses papéis são possibilidades aprendidas, e os pesos de uma cabeça, isoladamente, não provam por que o modelo fez determinada previsão.

\subsection{Da representação à geração}

Na tradução para o inglês, o codificador processa ``O banco aprovou o empréstimo'' e fornece os vetores de origem. O decodificador começa com um marcador inicial e usa a atenção cruzada para consultar essa representação. Uma camada linear transforma o vetor do decodificador em escores sobre o vocabulário. A softmax produz as probabilidades usadas para escolher o próximo token.

\begin{table}[htbp]
    \centering
    \begin{tabular}{>{\raggedright\arraybackslash}p{6.5cm} >{\raggedright\arraybackslash}p{6.5cm}}
        \toprule
        \textbf{Prefixo disponível} & \textbf{Próxima unidade ilustrativa} \\
        \midrule
        Marcador inicial & The \\
        \addlinespace
        The & bank \\
        \addlinespace
        The bank & approved \\
        \addlinespace
        The bank approved & the \\
        \addlinespace
        The bank approved the & loan \\
        \bottomrule
    \end{tabular}
\end{table}

A tabela usa palavras como unidades e apresenta uma sequência ilustrativa, sem simular as probabilidades de um modelo treinado. Em uma execução, a escolha depende da estratégia de decodificação. Cada token escolhido entra no prefixo e a geração continua até um marcador de fim ou outro critério de parada.

Durante o treino, a sequência alvo conhecida permite calcular previsões de várias posições em paralelo. O alvo é deslocado e a máscara causal bloqueia o acesso às posições à frente. Na geração, cada passo depende do prefixo já produzido. Modelos que usam apenas o decodificador também aplicam atenção causal, mas operam sobre o próprio contexto de entrada e de saída, sem atenção cruzada a um codificador separado.

\subsection{Vantagens e limites}

A autoatenção oferece um caminho direto entre posições distantes acessíveis e permite calcular as relações da sequência com operações matriciais. Em uma camada de atenção densa, porém, o número de pares de posições cresce com o quadrado do comprimento da sequência. Aumentar o contexto torna esse cálculo mais caro. Além disso, a representação contextual é uma combinação aprendida de vetores. A qualidade da resposta continua a depender dos dados, do treinamento e da estratégia de geração.

\section*{Referências}

\noindent [1] VASWANI, A. et al. \textbf{Attention Is All You Need}. NeurIPS, 2017. Disponível em: \url{https://arxiv.org/abs/1706.03762}.

\end{document}
